\section*{結言}
この日の特徴は、モデルそのものの性能競争だけでなく、AI開発の透明性、専門領域へのアクセス制御、複数エージェントを束ねる作業環境といった「運用の制度化」が前面に出た点にある。OpenAIのmisalignment開示枠とAnthropicのR\&D指標公開は異なる角度から、フロンティアAI開発の可視化を試みている。

同時に法務、生命科学、生体分子計算、ゲノム解析、3D空間生成、音声対話へと用途別の展開も続いた。未確認のベンチマークや将来製品の予告も残っており、翌日以降はDevDay周辺の公式発表と、今回公開された制度・製品の実運用情報が焦点となる。
